\section{Jordan Canonical Form}\label{sec:jordan-canonical-form}
Chapter~5 showed that when a split characteristic polynomial does not supply
enough ordinary eigenvectors, generalized eigenspaces recover the missing
directions, and one Jordan chain produces one Jordan block. The preceding
section proved
\[
V=\bigoplus_\lambda G_\lambda(T).
\]
The remaining question is whether each generalized eigenspace contains enough
compatible chains. The PID structure theorem answers it exactly.

\begin{theorem}[Jordan canonical form]\label{thm:jordan-canonical-form}
If \(\chi_T\) splits over \(F\), then some basis \(\beta\) gives
\[
[T]_\beta=\bigoplus_{\lambda,j}J_{k_{\lambda,j}}(\lambda).
\]
The multiset of blocks is unique up to permutation.
\end{theorem}
\begin{proof}
Every irreducible factor occurring in \(V_T\) is linear, so the
elementary-divisor decomposition is
\[
V_T\cong\bigoplus_{\lambda,j}
F[x]/((x-\lambda)^{k_{\lambda,j}}).
\]
The \(\lambda\)-summands form \(G_\lambda(T)\). By
Proposition~\ref{prop:jordan-chain-cyclic-module}, each cyclic summand
supplies a Jordan-chain basis and one block
\(J_{k_{\lambda,j}}(\lambda)\). Concatenating these bases proves existence.
Uniqueness of the elementary divisors proves uniqueness of the block
multiset, since \((x-\lambda)^k\) corresponds to exactly one
\(J_k(\lambda)\).
\end{proof}

\begin{corollary}[Kernel growth and block counting]
\label{cor:jordan-kernel-growth}
For \(K_j=\ker(T-\lambda I)^j\), \(K_0=0\), and
\(c_j=\dim K_j-\dim K_{j-1}\),
\[
\begin{aligned}
c_j&=\#\{\lambda\text{-blocks of size at least }j\},\\
c_j-c_{j+1}&=\#\{\lambda\text{-blocks of size exactly }j\}.
\end{aligned}
\]
Consequently the number of \(\lambda\)-blocks is
\(\dim E_\lambda\), and their total size is
\[
\dim G_\lambda(T)=d_\lambda,
\]
the algebraic multiplicity. The largest block size is the exponent of
\(x-\lambda\) in \(m_T\).
\end{corollary}
\begin{proof}
A block of size \(k\) contributes \(\min(j,k)\) to \(\dim K_j\), hence one
to \(c_j\) exactly when \(k\geq j\). Blocks for other eigenvalues contribute
zero. Summing proves the formulas. The remaining statements follow at
\(j=1\), by summing all block sizes, and by taking the least common multiple
of the elementary divisors.
\end{proof}

Since \(E_\lambda\subseteq G_\lambda(T)\) and
\(\dim G_\lambda(T)=d_\lambda\),
\[
E_\lambda=G_\lambda(T)\quad\Longleftrightarrow\quad
\dim E_\lambda=d_\lambda.
\]
Thus no generalized vectors beyond ordinary eigenvectors are needed exactly
when the eigenspace accounts for the full algebraic multiplicity.

\begin{corollary}[Diagonalizability and the minimal polynomial]
\label{cor:diagonalizable-minimal-polynomial}
\(T\) is diagonalizable over \(F\) if and only if \(m_T\) is a product of
distinct linear factors over \(F\).
\end{corollary}
\begin{proof}
If \(m_T\) has distinct linear factors, every elementary divisor divides
\(m_T\), so every exponent and every Jordan-block size is \(1\). Conversely,
a diagonal operator has minimal polynomial equal to the product of the
distinct factors \(x-\lambda\) belonging to its eigenvalues.
\end{proof}

\begin{example}[What the second kernel detects]
If \(d_\lambda=4\) and \(\dim E_\lambda=2\), the two block patterns are
\(3+1\) and \(2+2\). They have
\[
\dim\ker(T-\lambda I)^2=3\quad\text{and}\quad4,
\]
respectively.
\end{example}

\section{Rational Canonical Form}\label{sec:rational-canonical-form}
Jordan form may fail over the base field. Over \(\mathbb R\),
\[
A=\begin{pmatrix}0&-1\\1&0\end{pmatrix},\qquad \chi_A=x^2+1,
\]
has no real eigenvalues. Yet \(v=e_1\) and \(Av=e_2\) form a basis, while
\(A^2v=-v\). When eigenvectors are unavailable, an orbit
\(v,Tv,T^2v,\ldots\) can replace them.

\descriptivesubsection{From an orbit to a companion block}
For \(0\ne v\in V\), the cyclic subspace
\[
W=F[x]v=\operatorname{span}\{v,Tv,T^2v,\ldots\}
\]
is \(T\)-invariant: applying \(T\) to a finite linear combination of orbit
vectors produces another such combination. Finite dimensionality forces a
dependence among \(v,Tv,T^2v,\ldots\). If \(m=\dim W\), the first dependence
occurs when \(T^mv\) is adjoined, so
\[
\beta=(v,Tv,\ldots,T^{m-1}v)
\]
is a basis of \(W\). The coefficient of \(T^mv\) in that first relation is
nonzero; dividing by it normalizes the relation to
\[
T^mv+a_{m-1}T^{m-1}v+\cdots+a_1Tv+a_0v=0,
\]
or, solving for the first dependent orbit vector,
\[
T^mv=-a_0v-a_1Tv-\cdots-a_{m-1}T^{m-1}v.
\]

Now read the action of \(T\) on \(\beta\):
\[
T(v)=Tv,\quad T(Tv)=T^2v,\quad\ldots,\quad
T(T^{m-2}v)=T^{m-1}v,
\]
while
\[
T(T^{m-1}v)
=-a_0v-a_1Tv-\cdots-a_{m-1}T^{m-1}v.
\]
The first formula gives column \(1\) with a \(1\) in row \(2\), column \(2\)
with a \(1\) in row \(3\), and so on through column \(m-1\), with a \(1\)
in row \(m\). The last formula gives the final coordinate column
\((-a_0,-a_1,\ldots,-a_{m-1})^{\mathsf T}\). Therefore
\[
[T|_W]_\beta=C_d=
\begin{pmatrix}
0&0&\cdots&0&-a_0\\
1&0&\cdots&0&-a_1\\
0&1&\cdots&0&-a_2\\
\vdots&&\ddots&\vdots&\vdots\\
0&0&\cdots&1&-a_{m-1}
\end{pmatrix},
\quad d=x^m+a_{m-1}x^{m-1}+\cdots+a_0.
\]
The companion matrix is therefore not an artificial matrix recipe: it is
forced by the natural orbit basis of a cyclic operator. The first relation
also says that no polynomial of degree below \(m\) annihilates \(v\), while
\(d(T)v=0\). Hence
\[
\operatorname{Ann}(v)=(d),\qquad F[x]v\cong F[x]/(d).
\]
Hence
\[
\boxed{\text{cyclic invariant subspace}\longleftrightarrow
F[x]/(d)\longleftrightarrow C_d}.
\]
The parallel with the split case is concise:
\[
\text{Jordan chain}\longrightarrow\text{Jordan block},
\qquad
\text{cyclic orbit}\longrightarrow\text{companion block}.
\]

If \(V=F[x]v\), then \(\operatorname{Ann}(v)=\operatorname{Ann}(V_T)\), so
\(m_T=d\). Cayley--Hamilton gives \(m_T\mid\chi_T\), while
\[
\deg m_T=\deg d=\dim V=\deg\chi_T.
\]
Both polynomials are monic, hence
\[
m_T=\chi_T=d.
\]

\begin{theorem}[Characteristic polynomial of a companion matrix]
\label{thm:companion-characteristic}
\(\chi_{C_d}=d\).
\end{theorem}
\begin{proof}
Put \(D_m(x)=\det(xI-C_d)\). Explicitly,
\[
xI-C_d=
\begin{pmatrix}
x&0&\cdots&0&a_0\\
-1&x&\cdots&0&a_1\\
0&-1&\ddots&0&a_2\\
\vdots&\ddots&\ddots&x&\vdots\\
0&\cdots&0&-1&x+a_{m-1}
\end{pmatrix}.
\]
Expand along the first row. Only its first entry \(x\) and final entry
\(a_0\) are nonzero. Deleting the first row and first column leaves the
companion-type determinant
\[
D_{m-1}(x)
=x^{m-1}+a_{m-1}x^{m-2}+\cdots+a_2x+a_1.
\]
Thus the first contribution is \(xD_{m-1}(x)\). The minor belonging to
\(a_0\) is upper triangular with every diagonal entry equal to \(-1\), so
its determinant is \((-1)^{m-1}\). The cofactor sign is
\((-1)^{1+m}\), and
\[
(-1)^{1+m}(-1)^{m-1}=(-1)^{2m}=1.
\]
Hence the \(a_0\)-contribution is \(+a_0\), and
\[
D_m(x)=xD_{m-1}(x)+a_0.
\]
The base case is \(D_1(x)=x+a_0\). Induction therefore gives
\[
D_m(x)=x^m+a_{m-1}x^{m-1}+\cdots+a_1x+a_0=d(x).
\]
\end{proof}

What if no single cyclic vector generates all of \(V\)? The invariant-factor
decomposition supplies several cyclic pieces.

\begin{theorem}[Rational canonical form]\label{thm:rational-canonical-form}
Over every field, unique monic nonconstant invariant factors
\(d_1\mid\cdots\mid d_r\) give a basis in which
\[
[T]=C_{d_1}\oplus\cdots\oplus C_{d_r}.
\]
Moreover \(m_T=d_r\) and \(\chi_T=d_1\cdots d_r\).
\end{theorem}
\begin{proof}
The invariant-factor theorem gives
\[
V_T\cong F[x]/(d_1)\oplus\cdots\oplus F[x]/(d_r),
\qquad d_1\mid\cdots\mid d_r.
\]
Translating back to linear algebra gives
\[
V=W_1\oplus\cdots\oplus W_r,
\]
where each \(W_i\) is \(T\)-invariant and cyclic. Choose a cyclic generator
\(v_i\) and the orbit basis
\[
\beta_i=(v_i,Tv_i,\ldots,T^{\deg d_i-1}v_i).
\]
The preceding derivation gives
\([T|_{W_i}]_{\beta_i}=C_{d_i}\). Concatenating the \(\beta_i\) gives a
basis \(\beta\) of \(V\) with
\[
[T]_\beta=C_{d_1}\oplus\cdots\oplus C_{d_r}.
\]
This proves existence: invariant factors give cyclic subspaces, orbit bases,
and then companion blocks. Uniqueness of the invariant factors gives
uniqueness of the rational canonical form. Finally, the annihilator is
generated by \(\operatorname{lcm}(d_1,\ldots,d_r)=d_r\), while the companion
characteristic-polynomial theorem and direct sums give
\(\chi_T=d_1\cdots d_r\).
\end{proof}

\begin{theorem}[The \(xI-A\) computational bridge]
\label{thm:operator-cokernel-presentation}
If \(A=[T]_\beta\), then \(V_T\cong\operatorname{coker}(xI-A)\). If
\[
P(x)(xI-A)Q(x)=\operatorname{diag}(1,\ldots,1,d_1,\ldots,d_r)
\]
is Smith form over \(F[x]\), the nonunit entries are the invariant factors,
and
\[
\operatorname{RCF}(A)=\bigoplus_iC_{d_i}.
\]
\end{theorem}
\begin{proof}
Write \(\beta=(v_1,\ldots,v_n)\) and \(A=(a_{ij})\). Define an
\(F[x]\)-linear map
\[
\Phi:F[x]^n\longrightarrow V_T,\qquad \Phi(e_i)=v_i.
\]
It is surjective because the \(v_i\) form an \(F\)-basis and constants in
\(F[x]\) already produce every \(F\)-linear combination. Column \(j\) of
\(xI-A\) is
\[
xe_j-\sum_i a_{ij}e_i.
\]
Under \(\Phi\), it maps to
\[
x\cdot v_j-\sum_i a_{ij}v_i
=Tv_j-Tv_j=0.
\]
Therefore \(\operatorname{im}(xI-A)\subseteq\ker\Phi\), and \(\Phi\)
induces a surjective map
\[
\overline\Phi:\operatorname{coker}(xI-A)\longrightarrow V_T.
\]

For the inverse, send
\[
\sum_i c_iv_i\in V
\]
to the cokernel class of the constant vector
\((c_1,\ldots,c_n)^{\mathsf T}\). In the cokernel, the
column relations above say
\[
x[e_j]=\sum_i a_{ij}[e_i].
\]
This is exactly the coordinate formula for \(Tv_j\), so the proposed map
respects multiplication by \(x\), and therefore by every polynomial in
\(F[x]\). It is an \(F[x]\)-module map. The two composites fix the
generators \(v_i\) and \([e_i]\), respectively, so they are inverse and
\(V_T\cong\operatorname{coker}(xI-A)\).

Finally, multiplication on the left and right by the invertible matrices
\(P(x)\) and \(Q(x)\) does not change the cokernel up to isomorphism. The
Smith diagonal therefore presents
\[
\bigoplus_iF[x]/(d_i);
\]
unit entries present zero modules. Thus its monic nonunit diagonal entries
are exactly the invariant factors.
\end{proof}

\begin{example}[A complete rational-form calculation]
Over \(\mathbb R\), take
\[
A=\begin{pmatrix}0&-1&0\\1&0&0\\0&0&2\end{pmatrix},\qquad
xI-A=\begin{pmatrix}x&1&0\\-1&x&0\\0&0&x-2\end{pmatrix}.
\]
Let \(\Delta_1\) be the monic gcd of the entries, \(\Delta_2\) the monic gcd
of the \(2\times2\) minors, and \(\Delta_3=\det(xI-A)\). For Smith factors,
\[
d_1=\Delta_1,\qquad d_1d_2=\Delta_2,\qquad
d_1d_2d_3=\Delta_3.
\]
Here \(\Delta_1=1\) because \(xI-A\) contains the unit entry \(1\). Among
the \(2\times2\) minors are
\[
x-2,\qquad x(x-2),\qquad x^2+1.
\]
Any common divisor must divide both \(x-2\) and \(x^2+1\). These are
coprime in \(\mathbb R[x]\), since \(x-2\) does not divide \(x^2+1\):
evaluation at \(x=2\) gives \(5\ne0\). Hence \(\Delta_2=1\). Finally,
\[
\Delta_3=\det(xI-A)=(x^2+1)(x-2).
\]
Consequently \(d_1=d_2=1\) and
\(d_3=(x^2+1)(x-2)\). Thus the Smith form is
\(\operatorname{diag}(1,1,d)\), where
\(d=x^3-2x^2+x-2\), and
\[
\operatorname{RCF}(A)=C_d=
\begin{pmatrix}0&0&2\\1&0&-1\\0&1&2\end{pmatrix}.
\]
Over \(\mathbb C\), \(d=(x-i)(x+i)(x-2)\), so the Jordan form is
\(\operatorname{diag}(i,-i,2)\).
\end{example}

\descriptivesubsection{Using canonical forms to test similarity}
\[
\begin{aligned}
A\sim B
&\quad\Longleftrightarrow\quad
\operatorname{RCF}(A)=\operatorname{RCF}(B)\\
&\quad\Longleftrightarrow\quad
\text{\(A\) and \(B\) have the same invariant factors}
\end{aligned}
\]
over every field. When the characteristic polynomials split over \(F\),
\[
\begin{aligned}
A\sim B
&\quad\Longleftrightarrow\quad
\text{\(A\) and \(B\) have the same Jordan-block multiset}\\
&\quad\Longleftrightarrow\quad
\text{\(A\) and \(B\) have the same elementary divisors}.
\end{aligned}
\]
Efficient tests proceed through
\[
\chi\longrightarrow m\longrightarrow\dim E_\lambda
\longrightarrow\dim\ker(A-\lambda I)^k
\longrightarrow\text{full block data}.
\]
A mismatch at any stage proves nonsimilarity. Matching characteristic
polynomials is not sufficient: \(I_2\) and
\(\begin{pmatrix}1&1\\0&1\end{pmatrix}\) have the same characteristic
polynomial but different minimal polynomials. Even characteristic and minimal
polynomials together need not determine all blocks in higher dimensions.
For example, the nilpotent block patterns
\[
3+3\qquad\text{and}\qquad3+2+1
\]
both have characteristic polynomial \(x^6\) and minimal polynomial \(x^3\),
but they are not similar: their kernel dimensions are \(2\) and \(3\),
respectively. Solving \(B=P^{-1}AP\) directly is usually a poor first
strategy; construct \(P\) only when an explicit basis change is asked.

\descriptivesubsection{Jordan form versus rational form}
\[
\begin{array}{c|c}
\text{Jordan form}&\text{Rational form}\\ \hline
\chi_T\text{ splits}&\text{every field}\\
\text{elementary divisors}&\text{invariant factors}\\
F[x]/((x-\lambda)^k)&F[x]/(d_i)\\
\text{Jordan-chain basis}&\text{cyclic orbit basis}\\
J_k(\lambda)&C_{d_i}
\end{array}
\]
These are the matrix manifestations of the two PID decompositions. Factoring
\(d\) into coprime prime powers and applying CRT decomposes \(F[x]/(d)\)
into \(F[x]/(p^e)\). When \(p=x-\lambda\), those pieces become Jordan
blocks. Rational form is field-independent classification; Jordan form is
the split-polynomial refinement.
